\clearpage
\fancyhead[L]{\small CERES ELEMENTE XII.}
\hypertarget{native-304-lower}{}
\begin{center}CERES ELEMENTE XII.
\end{center}

\begin{center}
\fontsize{10}{12}\selectfont\begin{tabular}{c}
\hline
Göttingische gelehrte Anzeigen. Stück 40. Seite 393..395. 10. März 1808. \\
\hline
\end{tabular}
\end{center}

Die Ergebnisse der vorjährigen Beobachtungen der neuen Planeten Vesta, Pallas und Juno liegen der astronomischen Öffentlichkeit bereits vor. Die Ergebnisse zu Juno wurden erst vor Kurzem in diesen Blättern mitgeteilt. Es bleibt also nur noch Ceres, zu der Prof. \textsc{Gauss} seine Untersuchungen gerade abgeschlossen hat. Außer den Beobachtungen, die Prof. \textsc{Harding} im April und Mai 1807 an der hiesigen Sternwarte angestellt hat, und einigen Berliner Beobachtungen konnte Prof. \textsc{Gauss} noch eine schöne Reihe handschriftlich mitgeteilter Mailänder Beobachtungen heranziehen. Aus allen diesen Beobachtungen ergab sich für die letzte Opposition folgendes Resultat:


\begin{center}\fontsize{10}{12}\selectfont\setlength{\tabcolsep}{3pt}\renewcommand{\arraystretch}{1.15}\begin{tabular}{lr}
Zeit der Opposition 1807, Mai 3. .\;.\; & $4^{\text{u}}\,12^{\text{m}}\,44^{\text{s}}$\\
\multicolumn{2}{r}{Meridian von Göttingen}\\
Wahre Länge .\;.\;.\;.\;.\;.\;.\;,\;.\; & $222^\circ\;14'\;9''1$\\
Wahre geozentrische Breite .\;.\;.\;.\; & $10\;40\;15.9$ nordlich.\\
\end{tabular}\end{center}

\noindent Dies ist die \emph{fünfte} Opposition, die seit der Entdeckung der \emph{Ceres} beobachtet wurde.

\par\indent Die letzten Elemente, die auch in diesen Blättern (1806, S. 1946) [von \textsc{Harding}] mitgeteilt wurden, weichen von diesem Ergebnis nur um 3 Minuten in der Länge und nur um 13 Sekunden in der heliozentrischen Breite ab. So klein dieser Unterschied auch ist, so eindeutig beweist er doch, dass die Störungen, auf die sich Prof. \textsc{Gauss} bisher beschränkt hat, nicht mehr ausreichen, um alle vorliegenden Beobachtungen genau wiederzugeben. Bisher wurden nämlich nur die von Jupiter verursachten Störungen berücksichtigt, und auch dabei blieben jene ausgeschlossen, die von höheren Potenzen der Exzentrizitäten abhängen. Um die bisherigen und die künftigen Beobachtungen genau miteinander in Einklang zu bringen, wird man daher vor allem diese zuletzt genannten Störungen nachträglich berücksichtigen müssen. Aus mehr als einem Grund wird es allerdings ratsam sein, für diese Arbeit noch eine oder ein paar neue Oppositionen abzuwarten. Diesmal hat sich Prof. \textsc{Gauss} also noch damit begnügt, die letzten Elemente unter Berücksichtigung allein derjenigen Störungen der Länge und des Radiusvektors zu verbessern, die im Märzheft der Monatlichen Correspondenz von 1803 in Tafeln zusammengestellt sind, wobei er die Störungen der Breite vollständiger berechnet hat (diese sind bisher noch nicht gedruckt worden). So werden die Differenzen möglichst auf alle fünf Oppositionen verteilt. Das Ergebnis dieser Untersuchung sind die folgenden neuen Elemente (in der Reihe Nr. XII), die sich, wie man sehen wird, nur sehr wenig von den letzten unterscheiden. Die Epochen gelten für den Meridian von Göttingen.

\Needspace{18\baselineskip}
\SourceLeaf{305}
\begin{center}\fontsize{10}{12}\selectfont\setlength{\tabcolsep}{3pt}\renewcommand{\arraystretch}{1.15}\begin{tabular}{lrr}
& Mittlere Länge & Perihel\\\hline
1802 & $155^\circ\;28'\;15''7$ & $146^\circ\;30'\;22''1$\\
1803 & $233\;38\;7.1$ & $32\;23.3$\\
1804 & $312\;0\;49.5$ & $34\;24.9$\\
1805 & $30\;10\;40.9$ & $36\;26.1$\\
1806 & $108\;20\;32.3$ & $38\;27.4$\\
1807 & $186\;30\;23.7$ & $40\;28.6$\\
1808 & $264\;53\;6.1$ & $42\;30.2$\\
\end{tabular}\end{center}

\begin{center}\fontsize{10}{12}\selectfont\setlength{\tabcolsep}{3pt}\renewcommand{\arraystretch}{1.15}\begin{tabular}{lr}
Tägliche mittlere Bewegung .\;.\;.\;.\;.\; & $770''9354$\\
Exzentrizität 1806 .\;.\;.\;.\;.\;.\;.\;.\;.\; & $0.0785251$\\
Jährliche Abnahme .\;.\;.\;.\;.\;.\;.\;.\; & $0.00000583$\\
Logarithmus der großen Halbachse .\;.\;.\; & $0.4420439$\\
Aufsteigender Knoten 1806 .\;.\;.\;.\; & $80^\circ\;53'\;24''6$\\
Jährliche Bewegung .\;.\;.\;.\;.\;.\;.\;.\; & $+1''48$\\
Neigung der Bahn 1806 .\;.\;.\;.\;.\;.\; & $10^\circ\;37'\;33''2$\\
Jährliche Abnahme .\;.\;.\;.\;.\;.\;.\;.\; & $-0''44$\\
\end{tabular}\end{center}

\par\indent Mit den fünf bisherigen Oppositionen stimmen diese Elemente folgendermaßen überein:

\begin{center}\fontsize{9}{11}\selectfont\setlength{\tabcolsep}{3pt}\renewcommand{\arraystretch}{1.15}\begin{tabular}{lrrrr}
& \multicolumn{2}{c}{Unterschied} & \multicolumn{2}{c}{\shortstack{Nach den\\XIII. Elementen}}\\
& \shortstack{heliozentr. Länge\\in der Bahn} & \shortstack{heliozentrische\\Breite} & & \\\hline
1802 & $-24''1$ & $+1''2$ & $-0''5$ & $-1''9$\\
1803 & $+40.9$ & $-9.6$ & $+22.1$ & $-5.3$\\
1804 & $-21.0$ & $-3.1$ & $-34.9$ & $-1.2$\\
1806 & $+7.8$ & $0$ & $+26.7$ & $-3.6$\\
1807 & $-9.5$ & $-11.3$ & $-26.9$ & $-11.7$\\
1808 & $+62.7$ & & $+20.8$ & $-6.0$\\
\end{tabular}\end{center}

\noindent [Den Vergleich mit der Opposition von 1808 und die Unterschiede für die Elemente Nr. XIII hat \textsc{Gauss} handschriftlich hinzugefügt. \textsc{Schering}.]

